\chapter{Prise en main}
\section{Installation}
\begin{itemize}
  \item \textbf{Claude Code} : \cmd{claude --plugin-dir /chemin/vers/simetriik-forge}.
  \item \textbf{Codex et OpenCode} : lancer \cmd{./install.sh} (liens symboliques dans \cmd{~/.agents/skills}), puis redémarrer l'outil si les skills n'apparaissent pas.
\end{itemize}
Les skills s'invoquent par leur nom, par exemple \cmd{/simetriik-forge:forge-init} sous Claude Code.

\section{Démarrer un projet : \texttt{forge-init}}
Dans un dossier vide, lancez \cmd{forge-init}. Le skill :
\begin{enumerate}
  \item vérifie qu'il n'existe pas déjà de \cmd{docs/} ni de \cmd{AGENTS.md} (il n'écrase jamais) ;
  \item vous pose quatre questions, une à la fois : nom du projet, institution cliente, chef de projet, date de démarrage ;
  \item copie le dossier de conception (\cmd{docs/}, \cmd{api/}, \cmd{AGENTS.md}, \cmd{.github/}) et les scripts de garde-fou dans \cmd{scripts/} ;
  \item remplace les marqueurs \cmd{[NOM DU PROJET]}, \cmd{[INSTITUTION]}, \cmd{[NOM]}, \cmd{[JJ/MM/AAAA]} ;
  \item crée \cmd{CLAUDE.md} contenant \cmd{@AGENTS.md} ;
  \item annonce la suite : \cmd{forge-cadrage}.
\end{enumerate}
\astuce{La stack technique n'est \emph{pas} demandée à l'initialisation : c'est une décision d'architecture, prise à l'étape~3 par un ADR, à partir du cadrage et du domaine.}

\section{Structure du dossier de conception}
\begin{lstlisting}
docs/00-cadrage/        vision, impact map
docs/01-domaine/        event storming, glossaire, context map
docs/02-specs/          SPEC-NNN.md, features/*.feature, exigences non fonctionnelles
docs/03-architecture/   adr/, c4/workspace.dsl, fitness-functions.md
docs/04-livrables-client/
docs/05-recette/        plan-recette.md
docs/06-production/     deploiement, runbook, rollback, monitoring
api/openapi.yaml        contrat d'API
AGENTS.md               consignes pour les assistants de code
\end{lstlisting}

\section{Suivre sa progression}
Deux outils, tous deux en lecture seule :
\begin{itemize}
  \item \cmd{forge-statut} : tableau de bord complet (état de chaque étape et de chaque spec, garde-fous en échec) et \emph{une seule} prochaine action recommandée.
  \item \cmd{scripts/forge-progress.sh} : barre de progression sur les sept étapes, affichée à la fin de chaque skill, par exemple
\par\smallskip{\small\cmd{Progression [███░░░░] 3/7 · étape en cours : Architecture · suite : forge-contrat}}\par
\end{itemize}
Une étape compte comme faite lorsque son livrable est au statut \emph{Validée} (ADR : \emph{Acceptée}).
Pour l'architecture, la barre exige au moins deux ADR acceptés (l'ADR-0001 du monolithe modulaire et l'ADR de stack).

\section{Les statuts}
Brouillon $\rightarrow$ En revue $\rightarrow$ \textbf{Validée}. Seul un humain passe un livrable à \emph{Validée} (en indiquant son nom et la date).
\source{skills/forge-init, forge-statut}
